\section{Related work}
\label{sec:related}

\subsection{Normalization as a statistical interface}

Batch normalization \citep{ioffe2015bn} is usually discussed for what it does
during training. At inference it is a different object: a fixed affine map
whose slope and offset were estimated on the training cohort, so that the
representation of a test recording carries a term determined by how far that
recording's statistics lie from the cohort's. Group normalization
\citep{wu2018group}, instance normalization \citep{ulyanov2016instance} and
layer normalization \citep{ba2016layer} remove the batch from the computation
and make the output of the layer depend only on the example it is applied to.
\citet{frankle2021bn} show how much of a trained network's function is
reproduced by its normalization parameters alone, and \citet{nado2020}
evaluate prediction-time batch normalization under covariate shift; both
locate a large part of the model's behavior in statistics rather than in
filters.

Re-estimating the statistics on the test domain is adaptive batch
normalization \citep{li2018adabn}. In sleep staging the operation is
attractive because a night is a natural adaptation domain: its epochs are
exchangeable enough to estimate a mean and variance, and no label is needed.
What the operation cannot do is distinguish one sleep stage from another. It
estimates one mean per recording, so any patient-dependent component that
differs between stages survives it. This is the gap the present paper
addresses.

A related line conditions normalization on a label. Conditional batch
normalization and feature-wise transformations modulate the activations of a
network with a class or text embedding
\citep{devries2017modulating,dumoulin2018feature,perez2018film}; the
conditioning signal is known during training and supplied at inference. SCSN
conditions on the stage as well, but it does not learn a conditioning map and
it does not receive the stage: the stage is unknown at test time and is
replaced by the model's own soft posterior. The weights stay frozen, and the
only quantity estimated from the test recording is the shift itself.

The ingredients of the method are ordinary --- estimate a mean, subtract it,
index the estimate by a class --- so the question is what the paper contributes
beyond their assembly. The answer is the result that links them. The
decomposition of Section~\ref{sec:method} proves that no correction subtracting
one vector per recording can change the relative configuration of the
stage-conditional means, however well that vector is estimated
(Corollary~\ref{cor:no-marginal}); feature-wise modulation, in turn, is
supplied with a class taken from an annotation, whereas in staging the class is
unknown and the index has to be estimated from the model.
The two observations only become one contribution when the first shows that the
marginal route cannot work and the second shows what has to replace it. SCSN is
the normalizer that results: class-conditional statistics, conditioned on a
test-time estimate of the class rather than on a label. The experiments test
the derivation rather than a combination --- the interpolation between the
marginal and conditional endpoints, the matched-versus-mismatched control, and
the cross-cohort test in which a weak posterior should and does reduce the
method to its marginal special case.

\subsection{Test-time adaptation beyond marginal statistics}

Tent \citep{wang2021tent} updates the affine parameters of normalization
layers by minimizing the entropy of the posterior on the test recording;
test-time training \citep{sun2020ttt} uses a self-supervised objective. Both
change weights with an optimizer. CAFA \citep{cafa2023} aligns target features
to the source distribution with a class-aware objective, and UnMix-TNS
\citep{unmixtns2024} improves the estimation of normalization statistics when
test labels are temporally correlated. PSDNorm \citep{psdnorm2026} replaces
mean-and-variance normalization by a test-time alignment of the power spectral
density of each feature map and reports strong cross-dataset sleep staging;
the operation is input-domain and frequency-resolved, and it does not use the
predicted stage. FairAdaBN \citep{fairadabn2023} makes adaptive normalization
attribute-aware for fairness.

SCSN differs from this family in three respects. It is closed-form: no
optimizer, step size or number of steps is chosen at test time. It is
stage-conditional: the statistic it removes is indexed by the sleep stage, and
the index comes from the model's own posterior rather than from a label or a
frequency band. And its target is explicit: the class-conditional mean of the
representation, which is the quantity the classifier's decision boundaries were
fitted around. The interpolated form in Section~\ref{sec:method} contains the
post-hoc marginal correction as the endpoint $\lambda=0$, so the comparison is
a nested one rather than a comparison of unrelated methods.
Cross-domain sleep staging has also been addressed directly by test-time
adaptation \citep{guo2025tta}. The distinction from SCSN is the conditioning
variable: these methods adapt a shared target-domain or test-domain
transformation, whereas SCSN updates a separate first-moment correction for
each posterior-indexed sleep stage before the shared classifier is applied.

The family also has a blind spot that matters for the present paper. AdaBN,
Tent and their descendants operate on batch-normalization statistics or
affine normalization parameters. A group-normalized stager has no batch
statistic to replace; a bespoke entropy objective on its affine parameters
would be a different method from standard Tent and is not evaluated here.
Group normalization is an increasingly common choice precisely because it
removes dependence on batch composition, so the lack of a comparable
statistics-replacement correction is a practical gap. SCSN is defined on the
representation and therefore applies unchanged. The other correction compared
for such a checkpoint is post-hoc centering, which
removes one vector per recording; the experiments show that this removes the
overall subject probe while moving the within-stage probe by no more than a
third of what SCSN removes.

\subsection{Subject variability in automatic sleep staging}

Deep stagers have moved from single-cohort evaluation
\citep{supratak2017,phan2019seqsleepnet} to multi-cohort and cross-cohort
settings \citep{perslev2021,phan2022sleepformer}, and the gap between
within-cohort and cross-cohort accuracy has become an object of study.
\citet{nasiri2021subgrouping} show that population subgrouping improves staging
for under-represented patients. Adversarial subject classifiers
\citep{liao2020trifeaturenet} and disentangled architectures
\citep{s3dnet2026} change the weights so that the representation carries less
of the subject; both require retraining and, in the disentangled case, a
subject-invariance objective. \citet{cross2026} use subject-aware batch
normalization during training for cross-subject brain decoding. Normalization
has also been used as a cross-subject alignment device for EEG
\citep{katthi2021cca}, usually on the raw signal rather than inside the
network.
Recent Neurocomputing studies target subject generalization through few-shot
meta-learning \citep{metasth2026}, semi-supervised domain adaptation
\citep{sedaeeg2025}, and generalizable representation learning for EEG
\citep{transformeeg2026}. Like the methods above, these approaches alter
training or the learned representation; the inference-time decomposition
studied here is complementary and can be applied to their frozen outputs.

These methods answer the question \emph{how do we train a model that does not
encode the subject}. The present paper answers a different question: given a
model that already encodes the subject, which part of the encoding is
removable at inference, and which part is not. The answer separates a
recording-wide offset from a stage-conditional signature; the first is removed
by marginal adaptation and the second by SCSN, and only the second requires a
stage posterior. A disentangled architecture can be combined with SCSN, and we
test SCSN on group-normalized checkpoints for exactly that reason.

\subsection{Uncertainty and patient-level reliability}

Sleep staging is scored with a set of stages, and disagreement on boundary
epochs is documented \citep{rosenberg2013}. Calibration measures whether a
predicted probability means what it says \citep{guo2017}; conformal prediction
produces a set with a finite-sample coverage guarantee
\citep{vovk2005,angelopoulos2023gentle}. Adaptive and regularized scores
tighten the sets \citep{romano2020,angelopoulos2021raps}, and the guarantee has
been extended to time series and distribution drift
\citep{gibbs2021,stankeviciute2021}.

The guarantee is marginal over the exchangeable unit. When the unit is the
patient, a stager can honour it on average while under-covering particular
patients, and a representation that separates patients concentrates the
failures. We therefore report per-patient coverage, its dispersion and the
worst-covered patient, and we ask whether the correction that removes the
patient signature also equalises the coverage. The question is not whether the
guarantee is restored --- it remains marginal --- but whether the mechanism
that concentrates the failures is the mechanism that carries the patient.
